%!TEX root = ../main.tex
\section{Introduction}
\label{sec:intro}

The COVID-19 pandemic highlighted the need for scalable, remote health-assessment methods that minimize physical contact.
Speech is an attractive candidate for this kind of assessment: it can be recorded using ubiquitous consumer devices, while respiratory illness can alter airflow, resonance, phonation, and speaking dynamics in ways that are reflected in the acoustic signal. 
These properties have motivated growing interest in speech as a digital biomarker for respiratory health~\cite{cummins2018speech}. 
However, before such systems can be considered for screening or monitoring, they must demonstrate that they detect transient health-related changes rather than stable properties of the speaker, recording session, or speaking task. 
The distinction is not academic: in the largest evaluation of this kind, audio classifiers that predicted SARS-CoV-2 infection with ROC-AUC 0.846 fell to 0.619 once recruitment-driven confounders were matched, and did not improve on a simple symptom checker~\cite{coppock2024audio}.

Automatic cold detection provides a useful setting in which to study this problem. 
The URTIC corpus used in the ComParE 2017 Cold sub-challenge contains 28,652 audio chunks from approximately 630 people, divided into three official speaker-independent partitions~\cite{schuller2017interspeech}. 
In Train and Development, only 37 of 210 participants have a cold. 
Moreover, the corpus is cross-sectional: each person is recorded once, in exactly one health state. 
The thousands of correlated chunks therefore do not constitute thousands of independent health observations. 
Because speaker identity is perfectly confounded with observed health state, a classifier can improve its apparent ability to recognize illness by learning who is speaking rather than what illness sounds like.

The three official partitions are speaker-disjoint, but that protection does not extend to model development. 
Architecture, feature, fusion-weight, and threshold choices are all made on internal splits constructed by the practitioner, and the student release of URTIC provides no speaker identities with which to make those splits speaker-disjoint. 
An internal split that separates chunks rather than people therefore rewards speaker recognition as much as cold detection. 
The risk was recognized at the time. In their write-up for the 2017 challenge, Huckvale and Beke warned that a classifier trained on 10,000 extracts from only 210 speakers, a minority of them cold, may simply become a speaker recognizer~\cite{huckvale2017cold}.

We ask three questions. 
First, how strongly do speaker proxies, threshold selection, partition choice, and Development reuse affect the reported number? 
Second, which cold-related signals are captured by compact handcrafted features, trajectory representations, and foundation-model/CNN systems, and which of those signals are complementary? 
Third, which modeling and de-confounding approaches survive these controls, and what do their failures teach us about future small-sample speech-health studies? 
Our contributions follow the same order: 
(i) an evaluation audit showing why a validated speaker embedding does not automatically make transferred cluster identities valid; 
(ii) five hidden-Test submissions spanning three substantially different model families; and 
(iii) a controlled account of positive, negative, and invalidated results, centered on methodological lessons that extend beyond this particular challenge.
